\scanpage{8}
les \(t\) exceptionnels étant ceux pour lesquels le second membre de \((\unsure{1|3})\) est ambigu, i.e. qui ne sont pas de mesure nulle pour \(\unsure{\nu|t}\). Ainsi \(\varphi\) apparaît simplement comme fonction inverse de
\[
s\to\int_s^\infty d\unsure{\nu}(s).
\]

\deleted{\uuline{Définition}}
\uline{Corollaire} Soit \(f\) une fonction mesurable positive sur l'espace mesuré \(M\). Pour que \(f\) soit spectralement équivalente à une fonction décroissante \(\varphi\) sur \(R_*^+\), il faut et il suffit que pour tout \(a>0\), \(f^{-1}([a,+\infty[)\) soit intégrable. Alors \(\varphi_{\deleted{\unsure{f}}}\) est \unsure{unique|uniquement} \unsure{déterminée|déterminé}.

\deleted{\unsure{Définition}} On l'appelle le \uline{réarrangement décroissant} de \(f\), noté \(\varphi_f\). Si \deleted{\(\unsure{f_*}\)} \(g\) est une autre fonction mesurable positive, sur un espace mesuré \(N\), satisfaisant aux conditions ci-dessus, alors \(f\) et \(g\) sont \deleted{\unsure{spectralement}} équivalentes si et seulement si elles ont même réarrangement décroissant.
\deleted{\unsure{Si} \(\unsure{E}\) \unsure{est un} \illegible{3}}\inserted{\(\varphi_f(t)\) est défini par}
\[
(4)\qquad
\varphi_f^-(t)=\inf_{\substack{E\subset M\\m(E)\unsure{<|\leq}t}}\|(1-E)f\|_\infty
\qquad
\varphi_f^+(t)=\inf_{\substack{E\subset M\\m(E)\leq t}}\|(1-E)f\|_\infty
\]
où \deleted{\(\unsure{\varphi_f(t)}\)} on pose
\[
\varphi^-(t)=\lim_{\substack{s\to t\\s<t}}\varphi(s),
\qquad\deleted{\unsure{\text{continuité à droite}}}\inserted{\varphi_f^+(t)=\lim_{\substack{s\to t\\s>t}}\varphi(s)}.
\]
\inserted{\unsure{Ainsi on identifie}} un ensemble mesurable \(E\) avec sa fonction caractéristique, ce qui donne un sens à (4). De (4) on tire les formules (pour des fonctions mesurables positives sur un même espace mesuré \(M\)) \deleted{\unsure{f et g étant} \illegible{5}} :
\[
(5)\qquad
\left\{\begin{aligned}
\varphi_f&\leq\varphi_g\quad\deleted{\unsure{\text{implique}}}\inserted{\text{si}}\quad\deleted{\unsure{\varphi_f\leq\varphi_g}}\quad f\leq g\\
\varphi_{\lambda f}&=\lambda\varphi_f\qquad\text{si }\lambda\geq0\\
\varphi_{f^\alpha}&=(\varphi_f)^\alpha\qquad\text{si }\alpha\unsure{>|\geq}0\\
\varphi_{f+g}&\leq\varphi_f+\varphi_g\\
\varphi_{fg}&\leq\varphi_f\varphi_g
\end{aligned}\right.
\]
\margin{\(\int f\,dm=\int s\,d\mu_f(s)\) d'où\par \(\int f\,dm=\int\varphi_f\,dt\)\par\(\unsure{N_p}(f)=\unsure{N_p}(\varphi_f)\) etc.}
\margin{plus généralement pour toute fonction \deleted{\unsure{monotone}}\inserted{\unsure{croissante}} \(\alpha\) sur \(R^+\), \deleted{\unsure{existe}}\par\(\varphi_{\alpha(f)}=\alpha(\varphi_f).\quad\Leftarrow\)\par\deleted{\unsure{positive dans le cas} \(f\) \illegible{4} \(\unsure{g=|f|}\)}\par\deleted{\(\varphi_{|f|}=|\varphi_f|\) \unsure{etc}}}

\deleted{Dans la première formule}, (si \(\varphi_g\) existe \deleted{\(\unsure{\varphi}\)}\inserted{et \(f\leq g,\ \varphi_f\)} existe aussi; si \(\varphi_f\) et \(\varphi_g\) existent, \inserted{\(\unsure{\varphi_{\alpha(f)}}\),}\(\varphi_{f+g}\) et \(\varphi_{fg}\) existent).

Soit \(\psi\) une fonction décroissante sur \(R_*^+\), admettant la mesure spectrale \(\unsure{\mu_\psi}\). Pour que

\mspage{9}
